% =============================================================
% LIMITATIONS
% =============================================================

The validation reported in this article is qualitative and limited
to a single scene and a single event (the PRODES 2022 window over
a subset of the Brazilian arc of deforestation). All annotations
were placed by a single expert; no inter-annotator agreement, no
comparison against dense supervision, and no held-out test set
are provided. A quantitative evaluation with multiple annotators
and against dense ground truth is the subject of a separate
manuscript in preparation. Nothing in this article should be read
as a scientific claim about deforestation detection accuracy; the
validation establishes that the tool and the protocol produce
usable output, not that the output is optimal.

The tool assumes that the two images of a pair share the same
pixel grid. Co-registration and resampling are upstream concerns
and are not addressed by iSAGE-CD. Similarly, only bi-temporal
change (one reference and one frame per pair) is supported;
multi-date time series would require an extension of the
pair-anchored label format described in
Section~\ref{sec:label_format}.

The interface is a local PyQt5 desktop application and was tested
primarily on Windows. Training and prediction scripts are supplied
separately from the annotator itself; users adopting iSAGE-CD for
a different task must plug in their own training loop and their
own inference script, both configured to produce the per-frame
prediction PNGs in the class-index format described in
Section~\ref{sec:iteration}.

Finally, labels in iSAGE-CD are anchored to a pair rather than to
an image, which means the same visual polygon may carry different
labels under different references (for example, a cleared area
may be \emph{new} against one before-image and \emph{old} against
another). This is a design feature for change detection but is a
semantic shift relative to conventional single-image annotation
and should be accounted for when training analysts on the tool.
